% !TEX root = ../main.tex
\chapter{Resumen}

La computación cuántica representa un paradigma emergente con potencial transformador para la ciencia de datos. Sin embargo, persiste una brecha significativa entre las promesas teóricas de los algoritmos cuánticos y su aplicabilidad práctica en entornos empresariales reales. Este trabajo propone un estudio comparativo riguroso entre algoritmos cuánticos —concretamente métodos de kernel cuántico y el algoritmo QAOA— y sus equivalentes clásicos —SVM, Random Forests y heurísticas de optimización— aplicados a problemas reales de clasificación supervisada y optimización combinatoria, utilizando datos procedentes de una empresa privada. Se evalúa el rendimiento mediante métricas de precisión, eficiencia computacional y escalabilidad. Adicionalmente, se diseña un prototipo conceptual de flujo de trabajo híbrido cuántico-clásico orientado a la integración práctica en pipelines productivos. Los resultados buscan aportar evidencia empírica sobre las condiciones reales de ventaja cuántica en ciencia de datos aplicada.

{\bf Palabras clave:} computación cuántica, ciencia de datos, ventaja cuántica, aprendizaje automático cuántico, QAOA